\documentclass[10pt,a4paper]{amsart}
\usepackage[margin=19mm]{geometry}
\usepackage{amsmath,amssymb,mathtools}
\usepackage[scr=rsfs,cal=euler]{mathalfa}
\usepackage{xfrac}
\usepackage[unicode,hidelinks,bookmarks=false]{hyperref}
\newcommand{\ie}{i.e.\ }
\newcommand{\cf}{cf.\ }
\newcommand{\resp}{respectively\ }
\newcommand{\Subsection}{\subsection}
\providecommand{\nobreakdash}{\nobreak\hbox{-}}
\setlength{\emergencystretch}{2em}
\multlinegap0pt
\usepackage{bm}
\usepackage{bbm}
\usepackage{dsfont}
\usepackage{xspace}
\usepackage{cancel}
\usepackage{mathtools}
\usepackage{cleveref}
\usepackage{amsthm}
\makeatletter
\@ifundefined{lemma}{\newtheorem{lemma}{Lemma}}{}
\@ifundefined{theorem}{\newtheorem{theorem}{Theorem}}{}
\makeatother
\newcommand{\A}{\mathcal{A}}
\newcommand{\F}{\mathcal{F}}
\newcommand{\N}{\mathbb{N}}
\newcommand{\PP}{\mathcal{P}}
\newcommand{\intr}{\mathrm{int}}
\newcommand{\rk}{\mathrm{rk}}
\renewcommand{\H}{\mathcal{H}}
\renewcommand{\SS}{\mathbb{S}}
\renewcommand{\S}{\mathcal{S}}
\title[A topological proof of Ky Fan's covering lemma]{A topological proof of Ky Fan's covering lemma}
\author{Bogdan Chornomaz}
\date{Source: arXiv:2507.22184v1 (CC BY 4.0)}
\begin{document}
\maketitle

\noindent\emph{Source and licence.} A topological proof of Ky Fan's covering lemma. Version arXiv:2507.22184v1 (CC BY 4.0), distributed under the Creative Commons Attribution 4.0 International licence, as recorded in the arXiv licence metadata for this exact version and shown on the abstract page https://arxiv.org/abs/2507.22184v1. Full rights notice: licenses/arxiv-2507.22184-CC-BY-4.0.txt.\par\smallskip

\noindent\textit{Editorial scope.} Counted unit: the Theorem whose source label is \texttt{th-kyfan-covers3} in the article (source0.tex lines 451--455, source label \texttt{th-kyfan-covers3}), delivered together with the article's own complete original proof of it (source0.tex lines 456--485, one \texttt{proof} environment of 30 lines). No mathematics is written, repaired or re-derived: statement and proof are the article's own text line for line. Required context retained uncounted: lem-Ky-Fan-generalized (lines 249--253); th-kyfan-covers2 (lines 275--285). Every definition, display and label of those lines is retained unchanged. Omitted from this package and not redistributed: 1--248, 254--274, 286--450, 486--506. Nothing inside a retained excerpt is omitted or altered: every retained body is delivered exactly as the article's lines are written, so no omission receipt of any kind accompanies this package. Citations: the retained excerpts carry no citation command at all, so no bibliography entry of the article is transcribed and no reference is redistributed. Reference adaptation: none was needed and none was made; every reference inside the delivered unit resolves inside it, so no unresolved reference or citation is produced. Printed slips kept literal: no printed word, symbol, hypothesis or number of the counted statement or of its proof was altered. Layout and class: the article's own class is \texttt{article} with packages including authblk, amsthm, amsmath, amssymb, bbold, hyperref, caption, cleveref, xcolor, tikz; the wrapper is the lane's pinned \texttt{amsart} editorial wrapper at 10pt on A4, which restates the theorem-like environments the retained text needs and adds the article's own macro definitions that the retained text needs (\texttt{\textbackslash A}, \texttt{\textbackslash F}, \texttt{\textbackslash H}, \texttt{\textbackslash N}, \texttt{\textbackslash PP}, \texttt{\textbackslash S}, \texttt{\textbackslash SS}, \texttt{\textbackslash intr}, \texttt{\textbackslash rk}), with extra wrapper packages (bm, bbm, dsfont, xspace, cancel, mathtools, cleveref). The delivered numbering is the unit's own. Assets: no retained span contains a figure environment, a graphics-inclusion command, a PDF-page inclusion, a table or a TikZ picture; no image file is copied into this package, the package declares no asset dependency and no image file is redistributed with it. Licence: version arXiv:2507.22184v1 of this article is distributed under the Creative Commons Attribution 4.0 International licence, as recorded in the arXiv licence metadata for this exact version and shown on the abstract page https://arxiv.org/abs/2507.22184v1. A full rights notice is delivered at licenses/arxiv-2507.22184-CC-BY-4.0.txt. Characters: every retained excerpt is pure ASCII, so the delivered engine, XeLaTeX with its default Latin Modern text font, reports no missing character.
\begin{lemma}\label{lem-Ky-Fan-generalized}
	For an $n$-cover $\F$, let $\rho\colon H^\intr_\F\rightarrow \PP$ be a monotone map into a poset $\PP$ such that $\rho(s) \neq \rho(-s)$ for all $s\in H^\intr_\F$. Then there is $u\in \SS^n$ such that $\rho\circ s(u)\in \PP$ has rank at least $(n + 1)/2$. %In particular, the rank of $\PP$ is at least $(n + 1)/2$.
	
	Furthermore, if $\PP = \H_{2, \infty}$, then there is $u\in \SS^n$ such that the rank of $\rho\circ s(u)\in \PP$ is at least $n$.
\end{lemma}

\begin{theorem}[Ky Fan's covering lemma for two orders]\label{th-kyfan-covers2}
	For an $n$-cover $\F$, let $\leq_1$ and $\leq_2$ be two linear orders on $\F$. Then there is $u\in \SS^n$ witnessing an alternating monotone pattern of length at least $m = \left((n+1)/2\right)^{1/2}$ for both of those orders. 
	
	That is, there is $u\in \SS^n$ and $F_1, \dots, F_m\in \F$ such that:
	\begin{itemize}
		\item Either $u\in (-1)^{i-1}F_i$, for $i=1, \dots, m$ or $u\in (-1)^{i}F_i$, for $i=1, \dots, m$;
		\item Either $F_1\leq_1 F_2 \leq_1 \dots \leq_1 F_m$, or $F_1\geq_1 F_2 \geq_1 \dots \geq_1 F_m$, and similarly for $\leq_2$.
	\end{itemize}
    
    Moreover, this is asymptotically sharp: there is an $n$-cover such that the largest alternating monotone pattern with respect to it has length at most $\lceil\sqrt{n+2}\rceil$.
\end{theorem}

\begin{theorem}[Ky Fan's covering lemma for multiple orders]\label{th-kyfan-covers3}
	For an $n$-cover $\F$ and $d\geq 1$, let $\leq_i$, for $i=1, \dots, d$ be $d$ linear orders on $\F$. Then there is $u\in \SS^n$ witnessing an alternating pattern of length at least $m = \left((n+1)/2\right)^{1/2^{d-1}}$, monotone with respect to all $d$ of these orders. 

    Moreover, this is asymptotically sharp: there is an $n$-cover such that the largest alternating monotone pattern with respect to it has length at most $\left\lceil(n+2)^{D}\right\rceil$, for $D = 1/2^{d-1}$.
\end{theorem}

\begin{proof}
	The proof follows the same pattern as the one in \Cref{th-kyfan-covers2}. We start by fixing a realizable sample $s\in \S_\F$.
	
	Let us define $R$ as a set of functions $r\colon [d]\rightarrow \{\pm 1\}$ such that $r(1) = 1$; note that $|R| = 2^{d-1}$. Now, for $r\in R$, we define $\iota_r \colon \sup(s) \rightarrow \N$ as the length of the maximal alternating, with respect to $s$, pattern, which is increasing in $\leq_{r(i)\cdot i}$ for all $i\in [d]$, and ends in $F$; here $\leq_{-i}$ is defined as the order, inverse to $\leq_i$.  That is, $\iota_r(F)$ is the maximal $k$ such that there is $F_1, \dots, F_k \in \sup(s)$ such that:
		\begin{itemize}
			\item $s(F_i) = -s(F_{i-1})$ for $i=1, \dots, k-1$;
			\item $F_i \leq_{r(j)\cdot j} F_{i+1}$ for all $i=1, \dots, k-1$ and $j\in [d]$;
			\item $F_k = F$.
		\end{itemize}
	Finally, we define $\iota \colon \sup (s) \rightarrow \N^R$ as $\iota(F) = \left(\iota_r(F)\right)_{r\in R}$. Similarly to how it was before, the following holds:
		\begin{equation}\label{obs-iota2}\tag{*}
		\text{\emph{For every $F, G \in \sup(s)$, if $s(F)\neq s(G)$, then $\iota(F)\neq \iota(G)$.}}
		\end{equation}	
		
	Again, similarly, we consider posets $\N^R$, $\A(\N^R)$, and $\PP^0$, where the latter is defined as a poset of tuples $(a, f)$, for $a\in \A(\N^R)$ an antichain in $\N^R$ and $f\colon a\rightarrow \{\pm 1\}$ arbitrary, where the order on $\PP^0$ is $(a, f)\leq (b, g)$ if $a<b$, or $a=b$ and $f=g$. Also, $\PP$ is  defined as $\PP^0 - \{(e, -), (e, +)\}$, where $e$ is all-$1$ function on $R$. And we define $\rho \colon \S_\F \rightarrow \PP$ as $\rho(s) = (a_s, f_s)$, where $a_s$ is the antichain of maximal elements in the image $\iota(s)\subseteq \N^R$, and $f_s \colon a_s\rightarrow \{\pm 1\}$ is such that $f_s(\iota(F)) = s(F)$ whenever $\iota(F)\in a_s$.
	
	Now, same as before, we apply \eqref{obs-iota2} to ensure that $\rho$ is well-defined and monotone.  Then, by \Cref{lem-Ky-Fan-generalized}, we get $u\in \SS^n$ such that for $w = (a,f) = \rho(s(u))$ we get $\rk_{\PP^0}(w) = \rk_{\A(\N^R)}(a) \geq (n+3)/2$. Then, for $r\in R$,  we define $h_r(a) = \max \{x(r)~|~x\in a\}$; recall that $a$ is an antichain in $\N^R$, and so $x\in a$ is an $R$-dimensional vector, so $x(r)$ makes sense. As before, $h_r$ is the size of the maximal alternating chain increasing in every $\leq_{r(i)\cdot i}$, for $i\in [d]$, which is witnessed by $u$. Then 
	 
	 \begin{align*}
	 	\rk_{\A(\N^R)} \left( \left\{ \left(h_r(a) \right)_{r\in R} \right\} \right)
	 		&= \prod_{r\in R} h_r(a) - 1 \geq  \rk_{\A(\N^R)}(a)
	 			\geq (n+3)/2,
	 \end{align*}
	 trivially implying 

	 \begin{align*}
	 		\max_{r\in R} h_r(a) &\geq \left((n+1)/2\right)^{1/|R|}
	 			= \left((n+1)/2\right)^{1/2^{d-1}}.
	 \end{align*}
\end{proof}

\end{document}
